\section{Background}\label{sec:background}
%%% ====================================================================

The introduction established five interconnected threads: (i)~the
Euler product as descent data connecting primes to spectral locations;
(ii)~the Lawvere--Frobenius double obstruction; (iii)~circumvention
strategies (Russell~\cite{russell03}, Priest~\cite{priest79}, Tarski~\cite{tarski52}); (iv)~convergence of
$\mathbb{F}_1$-geometry, string theory, and the Hilbert--P\'olya
conjecture on a common geometric need; (v)~tripartite factorization
as the proposed mechanism.  At the heart of the
$\mathbb{F}_1$-program lies a mirror duality---between moduli spaces
and function spaces. Manin's \emph{Numbers as Functions}~\cite{manin13}
demonstrates that ``the analogy
between~$\operatorname{Spec}\mathbb{Z}$ \ldots\ and algebraic curves
over finite fields'' can be made ``so elaborated and precise that one
could use a version of the technique of Andr\'e~Weil \ldots\ in order
to approach Riemann's conjecture''; Weil himself, in his 1940 letter
from prison, had proposed a ``Rosetta Stone'' connecting number fields,
function fields, and Riemann
surfaces~\cite{weil_letter40,milne15}.
To expose the structural content of this duality, we trace both
the geometric and the algebraic lineages of the modulus back to their
earliest documented sources, prior to the creation of complex numbers
and indeed prior to algebra, where the two traditions are still
visibly independent.  From this vantage point we then follow their
convergence in~$\mathbb{C}$, the dimensional limitations that
constrain algebraic extension, the decidability boundary and its
categorical formalization, the geometric precedent established by
Weil's proof~\cite{weil48}, and the physical mechanisms that realize dimensional
transcendence---identifying at each stage the structural principles
that will be mathematically developed in~\cref{sec:methods}--\cref{sec:categorical}.


%%% ====================================================================
%%% \S 2.1 — THE DUAL GENEALOGY OF THE MODULUS
%%% ====================================================================

\subsection{The dual genealogy of the modulus}\label{ssec:modulus}

The concept of \emph{modulus} enters mathematics through two
independent lineages---one geometric, the other algebraic---whose
convergence in the complex plane constitutes the primitive structure
that this paper exploits.

\subsubsection{The geometric modulus: from Vitruvius to Farish}\label{sec:vitruvian}

In modern mathematics, the moduli space~$\mathcal{M}_g$ classifies
Riemann surfaces of genus~$g$ up to conformal equivalence---each point
representing an entire equivalence class of complex structures.  The
concept is well established in algebraic
geometry, physics, and number theory alike.  Yet the term
itself---Latin \emph{modulus}, ``small measure''---is far older, and
its geometric meaning predates its algebraic one by nearly two
millennia.

Apollonius (\emph{Conics}, c.~200~BCE) organized the conic curves
through intersection properties, classifying an entire family of
geometric objects by their distinguishing parameters.
Marcus Vitruvius Pollio (\emph{De Architectura},
c.~25~BC, Book~IV, \S~3) applied this parametric logic to architecture,
defining the modulus as the semi-diameter of a
column at its base: the fundamental unit from which all architectural
proportions derive by rational multiples.  Vitruvius's system operates
through a precise geometric mechanism: a fixed origin (column center),
a radial measure extending from it, and proportions expressed as
multiples $k \cdot m$ of the basic module~$m$.  Column height,
intercolumniation, entablature, and pediment are all multiples or
fractions of the modulus (Doric $k = 7$ diameters, Ionic $k = 8.5$,
Corinthian $k \approx 9$--$10$).  Book~III connects this architectural
modulus with the proportions of the human body---foot $= 1/6$ height,
palm $= 1/24$, finger $= 1/96$.  Vitruvius's mechanism thus contains the triple
\begin{equation}\label{eq:triple}
  \{\text{origin},\;\text{modulus (radial magnitude)},\;\text{angle
  (direction)}\},
\end{equation}
which would---some eighteen centuries later---be fused with algebra in
the creation of the complex plane.

Pappus (\emph{Collection}, c.~320~CE) refined Apollonius's
classification of conics through the focus--directrix characterization,
reducing the entire family to a single eccentricity parameter.

The medieval period preserved Vitruvian knowledge through monastery
scriptoria.  Poggio Bracciolini recovered the most complete manuscript
of \emph{De Architectura} at St.~Gall Abbey (1416), bringing it
to Florence where Brunelleschi and Alberti studied it.

Filippo Brunelleschi (c.~1415) demonstrated systematic linear
perspective through his \emph{tavoletta} experiment---a painted panel of
the Florence Baptistery viewed through a peephole and mirror.  The
demonstration established that three-dimensional space projects onto a
two-dimensional surface according to geometric laws: the vanishing
point functions as origin, radial distance determines apparent size
($s \propto 1/d$), and angle determines position---a direct
manifestation of the Vitruvian mechanism~\cref{eq:triple} applied to
depth representation.  Brunelleschi's perspective is also the first
instance of a \emph{tower generated by a single modulus}: one
viewpoint produces an infinite recession of depth planes, each
smaller than the last in proportions that Renaissance painters
associated with harmonious ratios---a preformal realization,
prior to Desargues's formalization, of the principle that a fixed
generating unit can produce unbounded structural information.

Leon Battista Alberti codified these principles in \emph{Della pittura}~\cite{alberti}
(1436), treating painting as a ``window'' onto three-dimensional space,
synthesizing Roger Bacon's optical tradition, Alhazen's
\emph{perspectiva}, and Brunelleschi's demonstrations.  Sixteenth-century
architects refined Vitruvius's proportional system---whose smallest
operational unit was the sixth of a module
(\emph{De Arch.}~IV.3.4--6)---by subdividing the module into thirty
\emph{minutae}~\cite{harthicks98}, the same Latin term already
employed for the sixtieth of a degree.  The shared vocabulary between
angular and modular minutes formalized the fine subdivision of a
generating unit, one radial, the other directional---the two
components of the Vitruvian triple~\cref{eq:triple}.

Pacioli~\cite{pacioli1509}---whose \emph{De Divina Proportione} was
illustrated by Leonardo da Vinci---codified the Vitruvian principle of
a single generating module as a universal aesthetic law, with~$\golden$
as the privileged proportional constant.

The geometric modulus received systematic three-dimensional extension.
Girard Desargues (1639) formalized perspective through projective
geometry.  Gaspard Monge (1795)~\cite{monge1811geometrie} systematized three-dimensional
representation through three orthogonal views---plan, elevation,
profile---formalizing recovery of 3D information from 2D projections.
William Farish~\cite{farish22} (1822) published ``On Isometrical Perspective,''
providing the first mathematical rules for isometric drawing, in which
equal measures apply uniformly to height, width, and depth.  In
isometric projection, a cube viewed along the spatial diagonal $(1,1,1)$
projects as a regular hexagon in the plane, with three axes forming
$120^\circ$ angles.  Auguste Bravais~\cite{bravais48} (1848--1850) extended the lattice
mechanism to physical structure, demonstrating exactly 14~types of
three-dimensional lattices compatible with crystallographic symmetries.

These geometric systems---perspective, projective geometry, descriptive
geometry, isometric projection, crystallographic lattices---possess rich
structure: symmetry operations, metric relationships, transformation
rules.  They establish origins, employ radial measures, incorporate
angular components, and support rigid and scaling transformations.
Yet they carry \emph{no algebraic structure}: no field operations, no
multiplication, no division.  In some cases, as Frobenius would prove
(\cref{ssec:frobenius}), they \emph{cannot} carry algebraic structure.


\subsubsection{The algebraic modulus: Cardano to Gauss}

The parallel algebraic genealogy of the modulus traces from Cardano's
\emph{Ars Magna} (1545) and the resolution of cubic equations, through
Bombelli, Wallis, and Euler, toward the algebraic formalization of
$\sqrt{-1}$ as a legitimate mathematical entity.  However, $\mathbb{C}$
possesses a geometric aspect that developed simultaneously with this
algebraic formalization.  The geometric tradition---employing the
mechanism~\cref{eq:triple}---is what provided the term ``modulus'' to
complex analysis.

The first mathematical deployment of ``modulus'' in number theory
appears in Gauss's \emph{Disquisitiones Arithmeticae}~\cite{gauss1870disquisitiones} (1801), the last
major mathematical work composed in Latin.  Gauss defined modular
congruence: \emph{``Si numerus~$a$ numerorum $b$, $c$ differentiam
metitur, $b$ et $c$ secundum~$a$ congrui dicuntur \dots\ ipsum~$a$
modulum appellamus.''}  The arithmetic modulus~$m$ partitions
$\mathbb{Z}$ into equivalence classes: two numbers are ``the same''
modulo~$m$ if they differ by a multiple of~$m$.  While no explicit
evidence confirms Gauss's~\cite{gauss31} conscious borrowing from Vitruvius, the term
captures a shared conceptual role: a reference magnitude that organizes
a space.


\subsubsection{Argand's synthesis and Riemann's formalization}

Jean-Robert Argand~\cite{argand1874essai} (1806) unified the geometric and algebraic moduli.
His central contribution was the explicit separation of magnitude and
direction: every complex number is the product of a modulus (pure
magnitude) and a direction factor (pure angle):
$z = |z| \cdot (\cos\theta + i\sin\theta)$.  Argand established that
multiplication by~$i$ is equivalent to rotation by $90^\circ$ in the
plane.  No direct evidence survives that Argand consciously borrowed
the term from Vitruvius.  Yet the coincidence is
structurally dense: the name is identical (\emph{modulus}), the
components are identical (origin, radial magnitude, angle), and the
milieu is precisely the one in which the Vitruvian module and its
Renaissance refinement into \emph{minutae} had remained
professionally current for two centuries.  Argand was born in
Geneva---home to the world's first watchmakers' guild
(1601)---and was himself probably a practitioner of horology:
MacTutor identifies the author of the \emph{Essai} as ``probably
an expert technician in the clock
industry''~\cite{mactutor_argand}.%
\footnote{The traditional identification of the author as
Jean-Robert Argand (1768--1822), a Genevan bookkeeper in Paris, is
disputed.  Schubring~\cite{mactutor_argand} argues that
the biographical data---age, profession, lack of positive
evidence---are inconsistent with the technical expertise displayed in
the \emph{Essai}.  For the present argument, only the content of the
1806 work and its Genevan--horological milieu are relevant.}
In horology the dial itself is organized by the
triple~\cref{eq:triple}---center, radius (hands), angle
(minutes)---and the term \emph{minute} designates simultaneously a
temporal unit and an angular subdivision of the module: historically
and categorically, the Vitruvian mechanism~\cref{eq:triple}
transmitted through Renaissance horology.

The primitive structure of~$\mathbb{C}$ upon which complex
algebra was mounted is the triple~\cref{eq:triple}: the same mechanism
that Vitruvius used to organize columns, that Brunelleschi used to
project depth, and that Farish used to draw three dimensions
isometrically.  The decisive novelty was that Argand's synthesis endowed
this mechanism with algebraic field operations ($+$, $\times$, $\div$),
creating an object simultaneously geometric space and algebraic field.

Riemann~\cite{riemann1857} (1857, ``Theorie der Abel'schen Functionen'') generalized from
individual complex numbers to families of geometric objects.  He
introduced \emph{Modul} to designate the parameters that characterize
equivalence classes of compact Riemann surfaces, defining the moduli
space~$\mathcal{M}_g$ as the quotient parametrizing surfaces of
genus~$g$ under conformal isomorphism.  This concept---spaces whose
points represent entire equivalence classes of geometric
objects---constitutes one of the deepest formalizations of the modulus
idea: from Vitruvius's single column radius to Riemann's universal
parameter space classifying all tori.


%%% ====================================================================
%%% \S 2.2 — THE MATHEMATICAL STRUCTURE OF C
%%% ====================================================================

\subsection{The mathematical structure of
$\mathbb{C}$}\label{ssec:C-structure}

\subsubsection{Algebraic, topological, and geometric structure}

The complex numbers $\mathbb{C} = \mathbb{R}[i]/(i^2+1)$ possess
simultaneous algebraic (field), topological (metric space), and
geometric (plane) structure.  The algebraic extension from~$\mathbb{R}$
is generated by the single relation $i^2 = -1$.

Every $z \in \mathbb{C}$ admits polar decomposition
$z = r e^{i\theta}$ with $r = |z| \geq 0$ and
$\theta = \arg(z) \in [0, 2\pi)$.  Euler's formula~\cite{euler55}
$e^{i\theta} = \cos\theta + i\sin\theta$ encodes rotation by
angle~$\theta$.  The modulus is multiplicative:
$|z_1 z_2| = |z_1||z_2|$, a property inherited from the Pythagorean
theorem.

The complex plane admits three equivalent representations, each
emphasizing different structural aspects:
(i)~Euclidean: $\mathbb{C} \cong (\mathbb{R}^2, d_{\mathrm{eucl}})$
with metric $d(z_1, z_2) = |z_1 - z_2|$;
(ii)~Algebraic: $\mathbb{C} = \mathbb{R}[x]/(x^2+1)$, a quotient of the polynomial ring in the sense of Noether~\cite{noether21}, emphasizing field operations;
(iii)~Polar: $\mathbb{C} \cong \mathbb{R}_+ \times S^1$, separating
magnitude from phase.

\subsubsection{The modulus formula and geometric--algebraic duality}

For $z = x + iy$, the modulus satisfies
$|z|^2 = z \cdot \bar{z} = x^2 + y^2$, where
$\bar{z} = x - iy$ denotes complex conjugation.  This formula unifies
geometric distance with algebraic operations: the modulus is a
quantity that exists prior to and independent of the field
structure---it is the Euclidean distance from the origin, a purely
geometric object.  Yet it connects to algebra through conjugation,
which is the unique non-trivial automorphism of~$\mathbb{C}$
over~$\mathbb{R}$.

\subsubsection{Lattices, tori, and moduli spaces}

A \emph{lattice} in~$\mathbb{C}$ is a discrete subgroup
$\Lambda = \mathbb{Z}\omega_1 + \mathbb{Z}\omega_2$ with
$\omega_1/\omega_2 \notin \mathbb{R}$.  Every lattice inherits the
metric and topological structure of~$\mathbb{C}$---distances, angles,
symmetries are well defined---but the lattice itself carries no field
operations: one cannot divide lattice points or take square roots
within~$\Lambda$.  Lattices are geometric objects embedded
in~$\mathbb{C}$, not algebraic substructures of it.

Two lattices are distinguished by their rotational symmetry.
The \emph{Gaussian lattice} $\mathbb{Z}[i]$, generated by $1$
and~$i$, has $90^\circ$ rotational symmetry (the cyclic group
$\mathbb{Z}_4$, generated by multiplication by~$i$).  The
\emph{Eisenstein lattice} $\mathbb{Z}[\omega]$, where
$\omega = e^{2\pi i/3}$, has $60^\circ$ rotational symmetry (the
cyclic group~$\mathbb{Z}_6$, generated by multiplication
by~$-\omega^2 = e^{i\pi/3}$).
These are the \emph{only} lattices in~$\mathbb{C}$ with rotational
symmetry beyond the trivial $180^\circ$---a rigidity theorem that will
prove consequential for the construction.

The quotient $\mathbb{C}/\Lambda$ is a complex torus, and the
\emph{modular parameter} $\tau = \omega_2/\omega_1$ (with
$\mathrm{Im}(\tau) > 0$) characterizes the torus up to conformal
equivalence.  The space of all such parameters, modulo the action
of~$\mathrm{PSL}(2,\mathbb{Z})$, is the moduli
space~$\mathcal{M}_{\mathrm{lat}}$---Riemann's moduli space for
genus-one surfaces.


%%% ====================================================================
%%% \S 2.3 — THE FROBENIUS BARRIER
%%% ====================================================================

\subsection{The Frobenius barrier: algebraic impossibility versus
geometric viability}\label{ssec:frobenius}

William Rowan Hamilton devoted years to ``multiplying triples.''
According to his letters, his sons would ask each morning during
October 1843: ``Well, Papa, can you multiply triples?''  He was forced
to reply with a ``sad shake of the head: No, I can only add and subtract
them.''  The breakthrough came on October~16, 1843, along Dublin's Royal
Canal: three-dimensional multiplication was impossible, but
four-dimensional multiplication could work if commutativity were
abandoned.  Hamilton carved $i^2 = j^2 = k^2 = ijk = -1$ into
Brougham Bridge, founding the quaternions~$\mathbb{H}$---inherently
four-dimensional, with noncommutative multiplication~\cite{hamilton}.

Frobenius's theorem~\cite{frobenius1878} (1878) closed the question definitively: no
three-dimensional associative division algebra over~$\mathbb{R}$
exists (the complete classification and its spectral consequences
are reviewed in~\cref{sssec:frobenius}).  The temporal convergence is striking:
during the precise decades 1822--1878, geometry systematically
developed sophisticated three-dimensional representational
systems---Farish's isometric projection (1822), Bravais's spatial
lattices (1848)---each possessing rich structure yet requiring no
algebraic multiplication, while algebra encountered an
insurmountable barrier at dimension~three.

The conclusion is not that three-dimensional structure is impossible,
but that Frobenius constrains \emph{algebraic} structure---division
algebras, field extensions, internal multiplication---not
\emph{geometric} structure: metric spaces, normed spaces, topological
manifolds, lattices.  A three-dimensional metric
space is perfectly legitimate; what cannot exist is a
three-dimensional \emph{field}.  This divergence between algebraic
impossibility and geometric viability is the asymmetry that the
construction of~\cref{sec:methods} exploits.

\medskip
\noindent\textit{The binary inheritance and its consequences.}\quad
Twentieth-century mathematics and physics developed almost entirely
within this constraint.  Von~Neumann's quantum mechanics (1927--32)~\cite{vonneumann32}
placed states in complex Hilbert spaces~$\mathcal{H}$ over~$\mathbb{C}$;
Turing's universal machine (1936) operates on binary tape~$\{0,1\}$;
Shannon's information theory~\cite{shannon48} (1948) measures entropy in bits;
Boolean algebra structures logic through binary truth values.
Where three-dimensional structure appears---$\mathrm{SU}(3)$ color
symmetry, spatial coordinates in relativity, three-generation
structure in particle physics---it manifests \emph{geometrically}
(rotation groups, Lie algebras acting on complex vector spaces), not
through internal multiplication within an extended number system.

The connection to the next section is direct: the binary structure
$f\colon T \times T \to Y$ that Frobenius confined us to is
\emph{precisely} the diagonal structure from which Lawvere's
paradoxes arise~\cite{lawvere69,yanofsky03}.  The construction $g(s) = \alpha(f(s,s))$ requires
binary auto-application---two copies of the same object composed
through the same operation.  Thus Frobenius's barrier and Lawvere's
obstruction are not independent constraints but two faces of a single
structural limitation: confinement to binary algebraic structure
simultaneously blocks dimensional extension and enables
self-referential paradox.  The question becomes: can the geometric
freedom that Frobenius leaves open also provide an escape from the
logical trap that binary structure creates?  This is what
Tarski's decidability boundary, developed next, addresses.


%%% ====================================================================
%%% \S 2.4 — THE DECIDABILITY BOUNDARY AND ITS CATEGORICAL FORMALIZATION
%%% ====================================================================

\subsection{The decidability boundary and its categorical
formalization}\label{ssec:decidability}

The introduction (\cref{ssec:circumvention}) identified three circumvention strategies
for the self-referential paradoxes that Lawvere--Yanofsky~\cite{lawvere69,yanofsky03} unified:
Russell's type stratification~\cite{russell08}, Priest's paraconsistency~\cite{priest79}, and Tarski's
decidable geometry~\cite{tarski52}.  The connection between these strategies was
characterized as a ``geometric escape''---a domain that is decidable
(Tarski), stratified (Russell), and therefore immune to the diagonal
constructions that generate paradox.  We develop here the deep
mathematical structure of this escape, focusing on the precise
mechanism by which self-referential capacity is lost and the
categorical formalization of the boundary at which this loss occurs.


\subsubsection{Tarski's decidability boundary: the arithmetic
of self-reference}\label{sssec:tarski}

Tarski~\cite{tarski52} (1952) showed that the first-order theory of
elementary algebra---equivalently, the theory of real closed
fields---admits elimination of quantifiers and is therefore
\emph{complete and decidable}: every well-formed sentence is either
provable or refutable, and there exists an algorithm that determines
which.  The same result extends to elementary Euclidean
geometry---points equipped with betweenness and equidistance---since
the geometric system is interpretable in the
algebraic one~\cite{tarskigivant99}.
By contrast, enriching the system with a predicate
$\mathrm{In}(x)$ denoting the property of being an integer renders it
undecidable: in Tarski's words, ``no decision method for the enlarged
system can be constructed''~\cite{tarski52}.  The boundary
lies at the capacity to define~$\mathbb{N}$ within the system.
Real closed fields possess both addition and multiplication, yet
remain decidable because they cannot define~$\mathbb{N}$ as a subset;
it is this structural incapacity to self-encode---not the absence of
any single operation---that places geometry outside the reach of
G\"odel's incompleteness.  Without G\"odel numbering, the system
cannot construct the diagonal function $f\colon T\times T\to Y$ that
Yanofsky~\cite{yanofsky03} identifies as the single categorical source
of the liar paradox, Russell's paradox, and Turing's halting problem.

The deeper insight concerns arithmetic itself.  G\"odel's
incompleteness requires \emph{both} addition and multiplication: the
G\"odel numbering $\ulcorner\phi\urcorner$ that encodes formulas as
numbers uses multiplicative encoding (prime factorization) composed
with additive sequencing (concatenation).  A system with addition
alone---Presburger arithmetic $(\mathbb{Z}, +)$---is
decidable~\cite{presburger29}.  The purely multiplicative theory of the
integers is likewise decidable, since the fundamental theorem of
arithmetic reduces it to the decidable theory of countably many
independent copies of the natural order\footnote{The decidability of
the multiplicative theory of~$\mathbb{N}$ follows from Skolem’s
analysis of multiplicative arithmetic~\cite{skolem31}; see also
Mostowski’s observation in Tarski~\cite[Note~5]{tarski52}
that elementary algebra, though lacking the general concept of an
integer, contains a decidable ``general arithmetic'' isomorphic to the multiplicative theory of natural numbers.}.  It is their
\emph{combination} $(\mathbb{Z}, +, \times)$ that generates
undecidability and self-reference.

This principle has a direct implication for the P\'olya operator
problem.  A construction that operates in a domain where addition has
been destroyed---where only multiplicative structure remains, as
in~$\Ztwenty$ (\cref{sec:categorical})---inhabits Tarski's decidable
fragment.  The Lawvere--Yanofsky diagonal (\cref{ssec:double-obstruction}) cannot form because
the arithmetic capacity for self-encoding is structurally absent.
This clarifies retroactively the strategies of Weil and Manin: Weil's
proof for function fields succeeded in the geometric (decidable)
setting; Manin's $\mathbb{F}_1$-program seeks to
place~$\mathrm{Spec}\,\mathbb{Z}$ in a purely multiplicative
framework where the same escape applies.


\subsubsection{Grothendieck's categorical formalization of the
boundary}\label{sssec:grothendieck}

Tarski identified \emph{where} the decidability boundary lies.
Grothendieck provided the \emph{arrow} that crosses it.

Grothendieck's reformulation of mathematics through categories,
functors, and topos theory~\cite{grothendieck1960elements,grothendieck_sga4} (1960s) introduced three structural
innovations directly relevant to the present work.

\emph{The functorial perspective.}  A moduli space should be
understood not as a set of isomorphism classes but as a \emph{functor}
$\mathcal{F}\colon \mathbf{Sch}^{\mathrm{op}} \to \Set$ assigning to
each scheme the set of families over it.  Geometric objects are
characterized by their relationships (morphisms) rather than their
internal structure (elements).  This shifts the focus from
objects---which are subject to Frobenius's dimensional
constraints---to morphisms between them---which are not.
Teichm\"uller's moduli spaces~\cite{teichmuller39} (\cref{ssec:C-structure}) received
their definitive formulation within this framework: the moduli
functor $\mathcal{M}_g$ classifying families of Riemann surfaces is
the precise categorical realization of Riemann's original concept.

\emph{The site and topos.}  Grothendieck's generalized topology
(SGA~\cite{grothendieck_sga4}) replaces open sets with covering families of morphisms.  A
topos---a category with the logical structure of $\Set$ but different
geometry---provides a framework where ``spaces'' can exist that have no
classical point-set realization.  This is precisely the framework that
$\mathbb{F}_1$-geometry would later require: spaces over the ``field
with one element'' (Tits~\cite{tits56}, Soul\'e~\cite{soule03}) that exist as functors rather than classical
geometric objects.  Borger's~\cite{borger09} identification of $\Lambda$-ring
structures as $\mathbb{F}_1$-descent data (\cref{ssec:convergence})
globalizes to what he describes as the big \'etale topos
over~$\mathbb{F}_1$---a Grothendieckian construction.

\emph{The forgetful functor: Tarski's boundary made categorical.}
Here the connection to Tarski becomes precise.  Categories of
mathematical objects are connected by functors that ``forget''
structure:
\[
  C^*\text{-}\mathbf{Alg}
    \xrightarrow{U_1} \Hilb \xrightarrow{U_2} \Met
    \xrightarrow{U_3} \Set.
\]
At each stage, algebraic structure is lost: $U_1$ forgets the
$C^*$-product, $U_2$ forgets linearity, $U_3$ forgets the metric.
The chain moves from rich algebraic structure (where
self-reference is possible and G\"odel's theorem applies) toward
pure geometry (where Tarski's decidability holds).

The forgetful functor \emph{implements} Tarski's boundary
categorically: it is the mechanism by which a system loses arithmetic
capacity---and with it, self-referential capacity---while potentially
preserving information content.  A functor is \emph{faithful} if it
is injective on morphisms: it forgets structure but not information.
The chain $C^*$-$\mathbf{Alg} \to \Hilb \to \Met \to \Set$ strips
away precisely the capacity for self-referential encoding (the
combination of addition and multiplication that G\"odel numbering
requires) while preserving the geometric and combinatorial data
(multiplicative structure, metric relations, set-theoretic
cardinalities) that carry spectral information.

This is the categorical content of the ``geometric escape'' identified
in the introduction (\cref{ssec:circumvention}): the domain where Lawvere's diagonal
\emph{cannot form} is not a prohibition but a \emph{categorical level}
reached by applying forgetful functors.  The P\'olya operator, if it
is to avoid the Lawvere--Yanofsky obstruction, must be constructed
at the $\Met$ or $\Set$ level---in the decidable fragment where the
arithmetic required for self-reference has been systematically
destroyed.


%%% ====================================================================
%%% \S 2.5 — PHYSICAL REALIZATIONS OF DIMENSIONAL TRANSCENDENCE
%%% ====================================================================

\subsection{Physical realizations of dimensional
transcendence}\label{ssec:physical}

The forgetful functor of~\cref{ssec:decidability} implements a
dimensional reduction---from arithmetic to geometry, from
self-referential to decidable.  Concrete realizations of analogous
reductions exist in three physical frameworks.


\subsubsection{String theory: compactification and the Virasoro
tower}\label{sssec:strings}

The two-dimensional worldsheet~$\Sigma$, parametrized
by~$\mathbb{C}$, embeds into $D$-dimensional target spacetime through
coordinate functions $X^\mu(\sigma,\tau)$.  The Polyakov action~\cite{polyakov81} is
conformally invariant, generating the infinite-dimensional Virasoro
algebra $\{L_n\}$ with central charge~$c$.  Information about
$D$-dimensional geometry is encoded in how the 2D surface curves
through the target---requiring no division algebra in dimension~3.

Dimensional compactification provides a complementary mechanism:
physical dimensions can exist as topological circles~$S^1$ without
manifesting as algebraic extensions.  A compactified dimension of
radius~$R$ ``disappears'' from the low-energy theory as $R \to 0$,
while its topological imprint (winding modes, Kaluza--Klein tower)
persists.

The Huerta--Schreiber tower~\cite{huerta2018m} (\cref{ssec:convergence}) derives M-theory's dimensional
hierarchy from successive maximal central extensions of the
superpoint, each level the unique consistent extension of the
previous.  This is Russell's tower principle realized in physics:
stratification through levels, with self-membership blocked across
levels.


\subsubsection{AdS/CFT: the holographic realization of the forgetful
functor}\label{sssec:adscft}

The AdS/CFT correspondence~\cite{maldacena98} (Maldacena, 1997) provides the most
concrete physical realization of the principle that Tarski identified
and Grothendieck formalized (\cref{sssec:grothendieck}).

A $(d{+}1)$-dimensional gravitational
theory in Anti-de~Sitter space is dual to a $d$-dimensional conformal
field theory on its boundary, through the group isomorphism
$\mathrm{SO}(d{+}1, 2) \cong \mathrm{Conf}(d)$.  Bulk fields
$\phi(x,z)$ project to boundary operators $\mathcal{O}(x)$ through
the radial limit $z \to 0$---the $(d+1)$th dimension is ``lost'' in
the projection while its information content is encoded in the
boundary's conformal structure.

The structural parallel with the forgetful functor chain is precise.
The forgetful functor $U_2\colon \Hilb \to \Met$ strips linearity
while preserving metric structure; the holographic projection strips
a spatial dimension while preserving conformal structure.  Both
implement the same principle: loss of structure (algebraic or
spatial) accompanied by preservation of information content.

The Bekenstein--Hawking entropy~\cite{bekenstein73,hawking75} $S = A/(4G_N)$ establishes that
information scales with boundary area, not bulk volume---the factor
$1/4 = 1/2^2$ in Planck units.  In AdS$_3$/CFT$_2$,
three-dimensional gravitational information is encoded on a
two-dimensional boundary, with the isometry group decomposing into
Virasoro algebras with central charge $c = 3\ell/(2G_N)$.

\begin{center}
\begin{tabular}{lcc}
\toprule
 & String Theory & AdS/CFT \\
\midrule
Structure & Embedding $\Sigma \subset$ spacetime
  & Projection Bulk $\to$ Boundary \\
Method & Coordinates $X^\mu(\sigma,\tau)$
  & Radial limit $z \to 0$ \\
Information flow & 2D $\to$ $D$ via functions
  & $(d{+}1) \leftrightarrow d$ via group isomorphism \\
\bottomrule
\end{tabular}
\end{center}

Both achieve dimensional extension without requiring division algebras
in dimension~3, operating within and through~$\mathbb{C}$.


\subsubsection{Connes--Douglas--Schwarz:
$\mathbb{F}_1 \cong D^1_{\mathrm{compact}}$}

The Connes--Douglas--Schwarz
equivalence~\cite{connes98}, introduced in~\cref{ssec:convergence}, connects toroidal
compactification to noncommutative geometry.  We make the structural
content explicit.

String theory's toroidal compactification in the limit $R \to 0$
produces a dimension that remains topologically present (maintaining
periodicity and winding modes) but becomes metrically collapsed.
Manin's $\mathbb{F}_1$-geometry, in the limit $q \to 1$, produces
structures that retain multiplicative periodicity (Frobenius lifts
$\psi_p$, roots of unity) but lose additive extension.  Both
limiting processes yield objects sharing the same formal properties:
degenerate $S^1$ fibers parametrizing discrete symmetries without
continuous extension.

This convergence has formal precedent: the Connes--Douglas--Schwarz
correspondence establishes that the noncommutative tori
$T(\theta)$ arising from toroidal compactification are precisely the
$C^*$-algebraic structures that Manin employed for his Real
Multiplication approach.  The identification
$\mathbb{F}_1 \cong D^1_{\text{compactified}}$ is therefore not an
analogy but a structural correspondence: dimensions with
multiplicative periodicity and no additive extension.

This characterizes the ``missing third dimension'' that
Frobenius (\cref{ssec:frobenius}) proves cannot exist as a division
algebra: it must be a dimension with discrete multiplicative scaling
($\golden$-mediated, as the construction of~\cref{sec:methods} will demonstrate) but
no field extension---a topological $S^1$ that does not correspond to
any algebraic element.


%%% ====================================================================
%%% \S 2.6 — THE GEOMETRIC PRECEDENT: WEIL'S PROOF
%%% ====================================================================

\subsection{The geometric precedent: Weil's proof for function
fields}\label{ssec:weil}

The introduction (\cref{ssec:double-obstruction}) established that the Riemann Hypothesis is a theorem for function fields---proved by Hasse~\cite{hasse34} (1934) for elliptic curves, Weil~\cite{weil40} (1940) for function fields, and extended to all curves by Weil~\cite{weil48} (1948)---and completed for arbitrary varieties by Deligne~\cite{deligne74} (1974) using the intersection form and the Lefschetz trace formula~\cite{milne15,oort14}. This geometric program was subsequently extended to $GL_n$ by Drinfeld~\cite{drinfeld77} and Lafforgue~\cite{lafforgue02} via the Langlands correspondence---whose arithmetic branch produced Wiles's proof~\cite{wiles95}---and the Arthur--Selberg trace formula~\cite{arthur89}. Deligne's proof for general varieties relies on the \'etale framework of Serre~\cite{serre55} and Grothendieck~\cite{grothendieck1960elements}. The aspiration to transport these methods from~$\mathbb{F}_q$ to~$\mathbb{Z}$ motivates the $\mathbb{F}_1$-program.  We develop here the \emph{technical
mechanism} of Weil's proof, because it is the template that the
categorical construction of~\cref{sec:categorical} reproduces.

\subsubsection{What Weil proved and how}

Let $C$ be a smooth projective curve of genus~$g$ over the finite
field~$\mathbb{F}_q$.  The Frobenius endomorphism
$\mathrm{Fr}\colon C \to C$ acts by $x \mapsto x^q$ on coordinates;
the number of $\mathbb{F}_{q^m}$-rational points is
$\#C(\mathbb{F}_{q^m}) = 1 - \sum_{i=1}^{2g}\alpha_i^m + q^m$,
where the $\alpha_i$ are the eigenvalues of Frobenius acting on
the first \'etale cohomology group $H^1_{\text{\'et}}(C, \mathbb{Q}_\ell)$.
The Riemann Hypothesis for~$C$ is the statement
$|\alpha_i| = q^{1/2}$ for all~$i$---the analogue of
$\mathrm{Re}(\rho) = 1/2$.

Weil's proof proceeds geometrically.  He constructs the arithmetic
surface $C \times C$ (the self-product of the curve) and establishes
an intersection form
$\langle\cdot,\cdot\rangle\colon \mathrm{Div}(C \times C) \times
\mathrm{Div}(C \times C) \to \mathbb{Z}$.
Two distinguished divisors---$\xi_0 = C \times \{P\}$ (a
``horizontal'' copy) and $\xi_1 = \{P\} \times C$ (a ``vertical''
copy)---satisfy $\langle\xi_0,\xi_0\rangle = \langle\xi_1,\xi_1\rangle = 0$
and $\langle\xi_0,\xi_1\rangle = 1$.
The Riemann--Roch theorem and Serre duality, applied to divisors
on~$C \times C$, produce an inequality
$\langle D, D\rangle \leq 2\,\langle D,\xi_0\rangle \cdot
\langle D,\xi_1\rangle$
for all divisors~$D$.  Taking $D$ to be the graph of the $m$-th
iterate of Frobenius, this inequality forces
$|\alpha_i| \leq q^{1/2}$; the functional equation provides the
complementary bound $|\alpha_i| \geq q^{1/2}$, yielding
$|\alpha_i| = q^{1/2}$~\cite{oort14,milne15}.

\subsubsection{The structural content}

Three features of Weil's proof merit emphasis, because they define
the requirements for any analogous construction over~$\mathbb{Z}$:

\begin{enumerate}[label=(\roman*)]
  \item \textbf{The arithmetic surface.}  The proof requires
    $C \times C$ as a geometric object with well-defined intersection
    theory.  For $\zeta(s)$ over~$\mathbb{Z}$, the analogue would be
    $\mathrm{Spec}(\mathbb{Z}) \times \mathrm{Spec}(\mathbb{Z})$---an
    ``arithmetic surface'' that does not exist in conventional algebraic
    geometry.  Manin's noncommutative tori~\cite{manin04} (\cref{ssec:convergence}) provide the
    geometric arena but lack the dynamical completion needed for the
    analogue of the Frobenius endomorphism.

  \item \textbf{The Frobenius endomorphism.}  The proof requires
    a canonical endomorphism whose eigenvalues are the quantities to
    be bounded.  Over~$\mathbb{F}_q$, this is
    $\mathrm{Fr}\colon x \mapsto x^q$---an \emph{intrinsic} operation
    determined by the arithmetic of the base field.  Borger's~\cite{borger09}
    $\Lambda$-ring structures (\cref{ssec:double-obstruction}) provide the characteristic-zero
    analogue: the composition law
    $\psi_p \circ \psi_q = \psi_{pq}$ mirrors the multiplicative
    structure of the Euler product itself.

  \item \textbf{The squeeze from intersection theory.}  The proof
    establishes $|\alpha_i| = q^{1/2}$ through \emph{two independent
    bounds}: an upper bound from intersection positivity and a
    lower bound from the functional equation.  Neither alone suffices;
    their conjunction produces equality.  This two-sided structure is
    the geometric precedent for the squeeze theorem that emerges in
    the categorical setting (\cref{sec:categorical}): two independent
    categorical bounds forcing $\mathrm{Re}(\rho) = 1/2$.
\end{enumerate}

The three features---arithmetic surface, canonical endomorphism,
two-sided squeeze---define the template that any analogous
construction over~$\mathbb{Z}$ must reproduce.  The categorical
construction of~\cref{sec:categorical} provides an explicit realization: two
independent bounds from the ternary structure forcing
$\mathrm{Re}(\rho) = 1/2$, mirroring Weil's squeeze from
intersection positivity and the functional equation.


%%% ====================================================================
%%% \S 2.7 — CONVERGENCE AND STRUCTURAL PRINCIPLES
%%% ====================================================================

\subsection{Convergence and structural
principles}\label{ssec:principles}

The developments of \cref{ssec:modulus}--\cref{ssec:weil}
exhibit a convergent architecture.  We make this convergence explicit
by identifying the structural parallels between domains and then
extracting the design principles for the construction of~\cref{sec:methods}.

\subsubsection{Cross-domain correspondences}

The parallel between Brunelleschi's perspective
(3D${}\to{}$2D encoding, depth lost as independent variable, encoded
as apparent size $s \propto 1/d$), Farish's isometric projection
($120^\circ$ angular encoding), and the holographic bulk--boundary
correspondence ($(d{+}1)\!\to\!d$ encoding through radial limit) is
not merely illustrative.  In all cases, a dimension is ``lost''
geometrically while its information content is preserved.
Grothendieck's forgetful functor (\cref{sssec:grothendieck})
formalizes this principle categorically; Tarski's decidability
theorem (\cref{sssec:tarski}) justifies it logically.

Russell's type tower~\cite{russell08} and the Huerta--Schreiber tower share the same
architectural pattern: hierarchical stratification through levels,
each uniquely determined by its predecessor, with self-membership
blocked across levels.  Borger's~\cite{borger09} Frobenius lifts
$\psi^k \circ \psi^\ell = \psi^{k\ell}$, surveyed by
Manin~\cite{manin13}, constitute the arithmetic realization.

Weil's proof (\cref{ssec:weil}) establishes the geometric
template: two independent bounds, obtained from intersection
theory and the functional equation, produce equality by squeeze.
The forgetful functor chain provides the categorical mechanism for
crossing from the arithmetic domain (where the Euler product lives)
to the geometric domain (where intersection-theoretic methods
apply)---precisely the transport that the $\mathbb{F}_1$-program
seeks to formalize~\cite{connes14}.

\begin{center}
\small
\begin{tabular}{@{}llll@{}}
\toprule
Domain & Strategy & Mechanism & Key insight \\
\midrule
Logic & Russell & Type hierarchy $U_n$
  & Tower prevents self-membership \\
Logic & Tarski & Decidable geometry
  & Arithmetic loss $\implies$ decidability \\
Mathematics & Grothendieck & Forgetful functor
  & Arrow crossing Tarski's boundary \\
Mathematics & Weil & Intersection on $C\!\times\!C$
  & Two bounds $\implies$ squeeze \\
Physics & Strings & Compactification
  & Dimensions as topology, not algebra \\
Physics & AdS/CFT & Holographic duality
  & $(d{+}1)\!\to\!d$ with info preserved \\
Arithmetic & Manin/$\mathbb{F}_1$ & $\Lambda$-rings
  & Numbers as functions on curves \\
\bottomrule
\end{tabular}
\end{center}


\subsubsection{Structural principles for the construction}

The preceding analysis establishes six
principles that any construction applied to the problem at hand
must satisfy.

\textbf{Principle 1 (Geometric modulus as generator).}
The Vitruvian triple
$\{\text{origin},\allowbreak\;\text{modulus},\allowbreak\;\text{angle}\}$
provides the primitive structure of~$\mathbb{C}$.  A construction
exploiting the geometric aspect of~$\mathbb{C}$ should use the
modulus as generating mechanism: a single constant from which
proportions derive by self-similar scaling (Vitruvius,
Brunelleschi, Argand).

\textbf{Principle 2 (Tower with constant modulus).}
Brunelleschi's perspective (the first projective system generated from
a single vanishing point), Russell's type
hierarchy~\cite{russell08,martinlof75}, the Virasoro tower,
and Borger's Frobenius lifts~\cite{borger09} all exhibit the same pattern: an infinite
sequence of levels generated by a fixed rule, with new information at
each level but a single governing modulus throughout.  Any tower
construction should have: (a)~a fixed generating ratio, (b)~preservation
of the fundamental modulus at every level, and (c)~structural
information growth.

\textbf{Principle 3 (Frobenius circumvention through geometry).}
No three-dimensional division algebra exists over~$\mathbb{R}$.  But
three-dimensional geometric structure exists abundantly (Farish,
Bravais).  Extension $\mathbb{C} \to$ 3D must operate as a scaling
modulus rather than a field extension.  String theory demonstrates this
concretely: compactified dimensions exist as topological~$S^1$ without
algebraic extension.  The identification
$\mathbb{F}_1 \cong D^1_{\text{compactified}}$ characterizes the
needed dimension as having multiplicative scaling but no additive
extension.

\textbf{Principle 4 (Lattice duality and the Euler product).}
Each lattice $\Lambda \subset \mathbb{C}$ defines a torus
$\mathbb{C}/\Lambda$---an elliptic curve.
The Eisenstein lattice ($\mathbb{Z}[\omega]$, $60^\circ$) carries
$\mathbb{Z}_6$ symmetry; its $\mathbb{Z}_3$ subgroup (the $120^\circ$
rotations $\{1,\omega,\omega^2\}$) provides the three-fold symmetry
that echoes the tripartite factorization of the Euler product
through~$\Ztwenty$.
The Gaussian lattice ($\mathbb{Z}[i]$, $90^\circ$) carries $\mathbb{Z}_4$
symmetry: the rotational structure native to~$\mathbb{C}$ as a field.
A construction connecting the Euler product's arithmetic
($\mathbb{Z}_3$) to the complex plane's geometry ($\mathbb{Z}_4$)
must mediate between these two lattices; the torus provides the arena
where both symmetries coexist, and the bridge coupling constant must
preserve metric relations while transforming angular structure from
hexagonal to square.  The construction in \S~3 employs~$\golden$
as coupling constant.

\textbf{Principle 5 (Holographic contraction: arithmetic to geometry).}
The forgetful functor chain
$C^*$-$\mathbf{Alg} \to \Hilb \to \Met \to \Set$
implements Tarski's decidability boundary categorically.
The functor destroys addition (and with it G\"odel numbering) while
preserving multiplication (and with it spectral structure).  Any
construction crossing from arithmetic to geometry must pass through
this chain, losing self-referential capacity while preserving the
multiplicative structure that carries the Euler product into the
ternary regime.

\textbf{Principle 6 (Squeeze from independent bounds).}
Weil's proof establishes $|\alpha_i| = q^{1/2}$ through two independent
bounds~\cite{oort14,milne15}.  Lawvere's theorem requires the diagonal morphism
$\Delta\colon A \to A \times A$ and point-surjective maps $A \to B^A$
to generate paradox.  If the categorical structure produces two
independent bounds on $\mathrm{Re}(\rho)$---neither alone
sufficient---their conjunction can force equality without the
self-referential machinery that a single direct proof would require.
